% GENERATED FILE. Edit canonical lessons, then run npm run generate:books.
% canonical-source-hash: fnv1a64:d9f0d04ae618e792
% canonical-lessons: FR-C462-orageux, FR-C462-foudre, FR-C462-grele, FR-C462-gronder, FR-C462-rafale

\chapter{L'orage gronde}
\label{ch:fr-l-orage-gronde}
\hlchaptermodality{462}

\begin{chapteropening}
\textbf{By the end of this chapter:} \emph{I can describe a storm: stormy weather, lightning, hail, rumbling and a gust of wind.}

Complete the last lesson of chapter 462: I can describe a storm: stormy weather, lightning, hail, rumbling and a gust of wind.
\end{chapteropening}

\section[orageux]{orageux — stormy, thundery}

\label{lesson:FR-C462-orageux}



\begin{quote}
Before the new one: say the French for the bright spell, then the French for to clear.
\end{quote}

\subsection*{You'll want to know: orageux}
\textbf{orageux} — \textquotedblleft{}stormy, thundery\textquotedblright{}.

\begin{quote}
Le temps est orageux depuis ce matin.
\end{quote}

\begin{quote}
\emph{The weather has been thundery since this morning.}
\end{quote}

Feminine \emph{orageuse}. Figurative: \emph{une discussion orageuse}.

\begin{cousinweb}
\emph{orageux} is built on \emph{orage} 'storm' (EtymDB lists it), from Old French \emph{ore} 'wind', which goes back to Latin \emph{aura} 'breeze'.
\end{cousinweb}

\begin{grammarlens}[title={What you've built}]
One more word for telling, explaining and describing at B1.
\end{grammarlens}

\subsection*{Your turn}
\begin{itemize}
\raggedright
  \item \emph{Say it:} \emph{orageux}
  \item \emph{Say it:} \emph{orageux}, once more
  \item \emph{Say it:} a sentence of your own with \emph{orageux}
  \item \emph{Recall:} say the French for the bright spell, then the French for to clear, then say \emph{orageux} again
\end{itemize}

\begin{tcolorbox}[breakable,colback=teal!4,colframe=teal!35!black,title={Before you move on}]
What does orageux mean? (\textquotedblleft{}Stormy, thundery\textquotedblright{}.) Say it once more.
\end{tcolorbox}
\section[la foudre]{la foudre — the lightning (that strikes)}

\label{lesson:FR-C462-foudre}



\begin{quote}
Before the new one: say the French for to clear, then the French for stormy.
\end{quote}

\subsection*{You'll want to know: la foudre}
\textbf{la foudre} — \textquotedblleft{}the lightning (that strikes)\textquotedblright{}.

\begin{quote}
La foudre est tombée sur un arbre du parc.
\end{quote}

\begin{quote}
\emph{Lightning struck a tree in the park.}
\end{quote}

\emph{La foudre} strikes; \emph{un éclair} is the flash you see.

\begin{cousinweb}
\emph{foudre} is inherited from Latin \emph{fulgur} 'lightning' (Lewis \& Short); EtymDB lists the word without a source.
\end{cousinweb}

\begin{grammarlens}[title={What you've built}]
One more word for telling, explaining and describing at B1.
\end{grammarlens}

\subsection*{Your turn}
\begin{itemize}
\raggedright
  \item \emph{Say it:} \emph{la foudre}
  \item \emph{Say it:} \emph{la foudre}, once more
  \item \emph{Say it:} a sentence of your own with \emph{la foudre}
  \item \emph{Recall:} say the French for to clear, then the French for stormy, then say \emph{la foudre} again
\end{itemize}

\begin{tcolorbox}[breakable,colback=teal!4,colframe=teal!35!black,title={Before you move on}]
What does la foudre mean? (\textquotedblleft{}The lightning (that strikes)\textquotedblright{}.) Say it once more.
\end{tcolorbox}
\section[la grêle]{la grêle — the hail}

\label{lesson:FR-C462-grele}



\begin{quote}
Before the new one: say the French for stormy, then the French for the lightning.
\end{quote}

\subsection*{You'll want to know: la grêle}
\textbf{la grêle} — \textquotedblleft{}the hail\textquotedblright{}.

\begin{quote}
La grêle a abîmé les voitures dans la rue.
\end{quote}

\begin{quote}
\emph{The hail damaged the cars in the street.}
\end{quote}

\emph{Il grêle} — it is hailing; \emph{un grêlon} — a hailstone.

\begin{cousinweb}
\emph{grêle} 'hail' is built on \emph{grêler} 'to hail', which is usually traced to a Frankish verb; the origin is not settled. It is not the same word as the adjective \emph{grêle} 'thin'.
\end{cousinweb}

\begin{grammarlens}[title={What you've built}]
One more word for telling, explaining and describing at B1.
\end{grammarlens}

\subsection*{Your turn}
\begin{itemize}
\raggedright
  \item \emph{Say it:} \emph{la grêle}
  \item \emph{Say it:} \emph{la grêle}, once more
  \item \emph{Say it:} a sentence of your own with \emph{la grêle}
  \item \emph{Recall:} say the French for stormy, then the French for the lightning, then say \emph{la grêle} again
\end{itemize}

\begin{tcolorbox}[breakable,colback=teal!4,colframe=teal!35!black,title={Before you move on}]
What does la grêle mean? (\textquotedblleft{}The hail\textquotedblright{}.) Say it once more.
\end{tcolorbox}
\section[gronder]{gronder — to rumble; to scold}

\label{lesson:FR-C462-gronder}



\begin{quote}
Before the new one: say the French for the lightning, then the French for the hail.
\end{quote}

\subsection*{You'll want to know: gronder}
\textbf{gronder} — \textquotedblleft{}to rumble; to scold\textquotedblright{}.

\begin{quote}
Le tonnerre gronde au loin.
\end{quote}

\begin{quote}
\emph{Thunder is rumbling in the distance.}
\end{quote}

\emph{Gronder un enfant} — to tell a child off.

\begin{cousinweb}
\emph{gronder} is inherited from Latin \emph{grundīre} 'to grunt' (EtymDB).
\end{cousinweb}

\begin{grammarlens}[title={What you've built}]
One more word for telling, explaining and describing at B1.
\end{grammarlens}

\subsection*{Your turn}
\begin{itemize}
\raggedright
  \item \emph{Say it:} \emph{gronder}
  \item \emph{Say it:} \emph{gronder}, once more
  \item \emph{Say it:} a sentence of your own with \emph{gronder}
  \item \emph{Recall:} say the French for the lightning, then the French for the hail, then say \emph{gronder} again
\end{itemize}

\begin{tcolorbox}[breakable,colback=teal!4,colframe=teal!35!black,title={Before you move on}]
What does gronder mean? (\textquotedblleft{}To rumble; to scold\textquotedblright{}.) Say it once more.
\end{tcolorbox}
\section[la rafale]{la rafale — the gust (of wind); the burst}

\label{lesson:FR-C462-rafale}



\begin{quote}
Before the new one: say the French for the hail, then the French for to rumble.
\end{quote}

\subsection*{You'll want to know: la rafale}
\textbf{la rafale} — \textquotedblleft{}the gust (of wind); the burst\textquotedblright{}.

\begin{quote}
Une rafale de vent a emporté son chapeau.
\end{quote}

\begin{quote}
\emph{A gust of wind blew his hat away.}
\end{quote}

\emph{Des rafales à cent kilomètres-heure}.

\begin{cousinweb}
The origin of \emph{rafale} is not settled. It is usually linked to Italian \emph{raffica} 'gust'; EtymDB lists it without a source.
\end{cousinweb}

\begin{grammarlens}[title={What you've built}]
One more word for telling, explaining and describing at B1.
\end{grammarlens}

\subsection*{Your turn}
\begin{itemize}
\raggedright
  \item \emph{Say it:} \emph{la rafale}
  \item \emph{Say it:} \emph{la rafale}, once more
  \item \emph{Say it:} a sentence of your own with \emph{la rafale}
  \item \emph{Recall:} say the French for the hail, then the French for to rumble, then say \emph{la rafale} again
\end{itemize}

\begin{tcolorbox}[breakable,colback=teal!4,colframe=teal!35!black,title={Before you move on}]
What does la rafale mean? (\textquotedblleft{}The gust (of wind); the burst\textquotedblright{}.) Say it once more.
\end{tcolorbox}
